\section*{अध्याय \ref{ch:b1:spherical-astronomy} --- गोलीय खगोल विज्ञान}

\begin{solution}{exo:b1:spherical-astronomy:1}
ध्रुव की ऊँचाई \omterm{def:b1:spherical-astronomy:geographic}{अक्षांश} के बराबर है: \ang{52}। कोई तारा कभी अस्त नहीं होता यदि
$\delta > \ang{90} - \ang{52} = \ang{38}$ हो, और कभी उदय नहीं होता यदि
$\delta < \ang{-38}$ हो।
\end{solution}

\begin{solution}{exo:b1:spherical-astronomy:2}
$\qty{6.75}{h}\times\ang{15} = \ang{101.25}$। $H = \mathrm{LST} - \alpha =
\qty{10}{h} - \qty{6}{h}\,\qty{45}{min} = \qty{3}{h}\,\qty{15}{min}$ है, यानी
\ang{48.75}। \omterm{def:b1:spherical-astronomy:equatorial}{घंटा कोण} धनात्मक है: तारा याम्योत्तर के पश्चिम में है, जिसे वह तीन और एक चौथाई
नाक्षत्र घंटे पहले पार कर चुका है।
\end{solution}

\begin{solution}{exo:b1:spherical-astronomy:3}
$h = \ang{90} - \varphi + \delta_\odot$: जून अयनांत पर \ang{73.44}, विषुवों पर \ang{50},
दिसंबर अयनांत पर \ang{26.56}।
\end{solution}

\begin{solution}{exo:b1:spherical-astronomy:4}
1900: 100 से विभाज्य है पर 400 से नहीं, अतः अधिवर्ष नहीं। 2000: 400 से विभाज्य, अधिवर्ष।
2100: अधिवर्ष नहीं। जूलियन वर्ष \omterm{def:b1:spherical-astronomy:tropical-year}{उष्णकटिबंधीय वर्ष} से $365.25 - 365.2422 = 0.0078$ दिन अधिक है,
यानी प्रति शताब्दी \num{0.78} दिन; ग्रेगोरी वर्ष $0.0003$ दिन, यानी प्रति शताब्दी \num{0.03} दिन।
\end{solution}

\begin{solution}{exo:b1:spherical-astronomy:5}
$\cos\Delta = \sin^2\ang{45} + \cos^2\ang{45}\cos\ang{90} = 0.5$, इसलिए
$\Delta = \ang{60}$ और दूरी $6371\times\pi/3 = \qty{6672}{km}$ है। समांतर के अनुदिश, त्रिज्या
$R\cos\ang{45}$ के \omterm{def:b1:spherical-astronomy:spherical-triangle}{लघु वृत्त} पर, दूरी
$6371\times\cos\ang{45}\times\pi/2 = \qty{7076}{km}$ है। \omterm{def:b1:spherical-astronomy:spherical-triangle}{महावृत्त} लगभग \qty{404}{km} छोटा है।
\end{solution}

\begin{solution}{exo:b1:spherical-astronomy:6}
$\sin h = \sin\ang{48}\sin\ang{20} + \cos\ang{48}\cos\ang{20}\cos\ang{30}
= 0.2542 + 0.5446 = 0.7987$, $h = \ang{53.0}$।
$\cos\mathrm{A} = (\sin\ang{20} - \sin\ang{48}\sin\ang{53.0})/(\cos\ang{48}
\cos\ang{53.0}) = (0.3420 - 0.5935)/0.4027 = -0.6246$, $\arccos = \ang{128.7}$;
तारा याम्योत्तर के पश्चिम में है ($H > 0$), इसलिए $\mathrm{A} = \ang{360} -
\ang{128.7} = \ang{231.3}$: दक्षिण-पश्चिम की ओर, क्षितिज से \ang{53} ऊपर।
\end{solution}

\begin{solution}{exo:b1:spherical-astronomy:7}
$\cos H_0 = -\tan\ang{60}\tan\ang{23.44} = -1.7321\times0.4335 = -0.7510$,
$H_0 = \ang{138.67}$, दिन की लंबाई $2\times138.67/15 = \qty{18.49}{h}$, लगभग
\qty{18}{h}~\qty{29}{min}, जैसा \ang{60} उत्तरी वक्र जून अयनांत पर दिखाता है।
\end{solution}

\begin{solution}{exo:b1:spherical-astronomy:8}
याम्योत्तर-पारण हर दिन \qty{3.93}{min} पहले होता है, इसलिए 30 दिनों में $30\times3.93 =
\qty{118}{min}$ पहले: तारा लगभग 20:02 पर पार करता है। दोपहर को पार करने के लिए उसे
\qty{10}{h} = \qty{600}{min} पहले खिसकना होगा, जिसमें $600/3.93 \approx 153$ दिन लगते हैं।
\end{solution}

\begin{solution}{exo:b1:spherical-astronomy:9}
(क) एक विषुव पर दोपहर का \omterm{def:b1:spherical-astronomy:horizontal}{उन्नतांश} \ang{45} है: $I = 1361\times\sin\ang{45}
= \qty{962}{W.m^{-2}}$। (ख) जून अयनांत पर $h = \ang{68.44}$ और
$(a/r)^2 = (149.598/152.100)^2 = 0.9674$:
$I = 1361\times0.9674\times\sin\ang{68.44} = \qty{1224}{W.m^{-2}}$। ऊँचा सूर्य बड़ी दूरी की
भरपाई से कहीं अधिक कर देता है।
\end{solution}

\begin{solution}{exo:b1:spherical-astronomy:10}
$\sin\delta_\odot = \sin\ang{23.44}\sin\ang{60} = 0.3445$,
$\delta_\odot = \ang{20.15}$। $\tan\alpha_\odot = \cos\ang{23.44}\tan\ang{60}
= 1.5890$, $\alpha_\odot = \ang{57.82}$ ($\lambda_\odot$ के समान चतुर्थांश में)।
\omterm{def:b1:spherical-astronomy:equatorial}{ऋजु आरोहण} \ang{2.18} पीछे है, अर्थात् $2.18\times4 \approx
\qty{8.7}{min}$ समय। \omterm{def:b1:spherical-astronomy:ecliptic}{क्रांतिवृत्त} के अनुदिश एक समान गति से चलने वाला सूर्य भी \omterm{def:b1:spherical-astronomy:equatorial}{ऋजु आरोहण} में
असमान गति करेगा, और वही मात्रा \omterm{def:b1:spherical-astronomy:equatorial}{घंटा कोण} तय करती है, और इसलिए धूपघड़ी: विषुव और अगले अयनांत
के बीच वह माध्य सूर्य के पीछे, और अयनांत तथा अगले विषुव के बीच आगे रहता है। यही समय के
समीकरण का तिर्यकता वाला पद है।
\end{solution}

\begin{solution}{exo:b1:spherical-astronomy:11}
उदय पर $h = 0$ है, अतः \cref{prop:b1:spherical-astronomy:transform} देता है
$\sin\delta = \cos\varphi\cos\mathrm{A}_0$; जून अयनांत पर $\delta = \varepsilon$, अतः
$\cos\mathrm{A}_0 = \sin\varepsilon/\cos\varphi$। $\mathrm{A}_0 = \ang{45}$ के लिए:
$\cos\varphi = \sin\ang{23.44}/\cos\ang{45} =
0.5626$, $\varphi = \ang{55.8}$ उ। $\ang{90} - \varepsilon =
\ang{66.56}$ से ऊपर, $\sin\varepsilon/\cos\varphi > 1$: उदय का कोई \omterm{def:b1:spherical-astronomy:horizontal}{दिगंश} नहीं होता, क्योंकि उस
दिन सूर्य अस्त नहीं होता (मध्यरात्रि का सूर्य)।
\end{solution}

\begin{solution}{exo:b1:spherical-astronomy:12}
(क) कोण $\theta$ वाला लून गोले का भाग $\theta/2\pi$ है:
$\frac{\theta}{2\pi}\times4\pi R^2 = 2\theta R^2$। (ख) लूनों के तीनों जोड़ों का कुल क्षेत्रफल
$2\times(2A + 2B + 2C)R^2$ है; वे गोले को एक बार ढकते हैं, सिवाय त्रिभुज और उसके विपरीत भाग के,
जिनमें से हर एक तीन बार ढका जाता है, यानी ज़रूरत से दोगुना:
$4(A + B + C)R^2 = 4\pi R^2 + 4S$, जहाँ $S$ त्रिभुज का क्षेत्रफल है। अतः
$S = R^2(A + B + C - \pi)$। (ग) तीन समकोण: $S = R^2(3\pi/2 - \pi) = \pi R^2/2$, जो गोले का
आठवाँ भाग (एक अष्टांश) है। आकृति का त्रिभुज: $89.9 + 64.9 + 67.4 - 180 = \ang{42.2}$, यानी
\num{0.736} रेडियन: $S = 0.736\,R^2$, गोले का \qty{5.9}{\percent}।
\end{solution}

\begin{solution}{pb:b1:spherical-astronomy:1}
\textbf{1.} \omterm{def:b1:spherical-astronomy:geographic}{सह-अक्षांश} $\ang{90} - (\ang{-24.63}) = \ang{114.63}$ और
$\ang{90} - \ang{51.48} = \ang{38.52}$; देशांतर-अंतर \ang{70.40}।

\textbf{2.} $\cos\Delta = \cos\ang{114.63}\cos\ang{38.52} +
\sin\ang{114.63}\sin\ang{38.52}\cos\ang{70.40} = -0.3260 + 0.1899 =
-0.1362$.

\textbf{3.} $\Delta = \ang{97.83}$; दूरी $6371\times1.7074 =
\qty{10878}{km}$, लगभग \qty{10900}{km}।

\textbf{4.}
\[
  \cos\mathrm{Az} = \frac{\cos\ang{38.52} - \cos\ang{114.63}\cos\ang{97.83}}
  {\sin\ang{114.63}\sin\ang{97.83}} = \frac{0.7823 - 0.0568}{0.9006} = 0.8056 ,
\]
$\mathrm{Az} = \ang{36.3}$, उत्तर-पूर्व की ओर।

\textbf{5.} $\sin\mathrm{Az} = \sin\ang{70.40}\sin\ang{38.52}/\sin\ang{97.83}
= 0.5922$, जो \ang{36.3} या \ang{143.7} देता है; प्रश्न 4 की धनात्मक कोज्या \ang{36.3} चुनती है।

\textbf{6.} नहीं। स्थिर दिशा हर याम्योत्तर को एक ही कोण पर काटती है और एक ध्रुव की ओर सर्पिल
बनाती है; \omterm{def:b1:spherical-astronomy:spherical-triangle}{महावृत्त} की दिशा लगातार बदलती है (वह केवल आरंभ में \ang{36.3} है)। विमान एक लंबा
मार्ग तय करता और कहीं और पहुँचता, जब तक वह अपनी दिशा ठीक न करता।

\textbf{7.} दक्षिण \omterm{def:b1:spherical-astronomy:celestial-sphere}{खगोलीय ध्रुव} का \omterm{def:b1:spherical-astronomy:horizontal}{उन्नतांश} \omterm{def:b1:spherical-astronomy:geographic}{अक्षांश} के परिमाण के बराबर है: \ang{24.63}।

\textbf{8.} $h = \ang{90} - |\ang{-24.63} - (\ang{-16.72})| = \ang{82.09}$;
$\delta > \varphi$ है, इसलिए सिरियस \omterm{def:b1:spherical-astronomy:celestial-sphere}{खमध्य} के \ang{7.91} उत्तर से गुज़रता है।

\textbf{9.} कोई तारा कभी अस्त नहीं होता यदि $\delta < -(\ang{90} - \ang{24.63}) =
\ang{-65.37}$। कैनोपस, \ang{-52.70} पर, परिध्रुवीय नहीं है: वह अस्त होता है।

\textbf{10.} $\ang{90} - |\ang{-24.63} + \ang{52.70}| = \ang{61.93}$, \omterm{def:b1:spherical-astronomy:celestial-sphere}{खमध्य} के दक्षिण में।

\textbf{11.} $\cos H_0 = -\tan(\ang{-24.63})\tan(\ang{-16.72}) =
-(0.4584)(0.3004) = -0.1377$; $H_0 = \ang{97.92}$। सिरियस ऊपर रहता है
$2\times97.92/15 = \qty{13.06}{h}$ \omterm{def:b1:spherical-astronomy:sidereal-time}{नाक्षत्र काल} तक, लगभग \qty{13}{h}।

\textbf{12.} $\cos\mathrm{A}_0 = \sin(\ang{-16.72})/\cos\ang{24.63} =
-0.3165$: सिरियस \omterm{def:b1:spherical-astronomy:horizontal}{दिगंश} \ang{108.5} पर उगता है, पूर्व के दक्षिण में।

\textbf{13.} जब $H = 0$, तब $\mathrm{LST} = \alpha = \qty{6}{h}~\qty{45}{min}~\qty{9}{s}$।

\textbf{14.} $\qty{86400}{s}\times365.2422/366.2422 = \qty{86164.09}{s} =
\qty{23}{h}~\qty{56}{min}~\qty{4}{s}$; याम्योत्तर-पारण हर रात
\qty{235.9}{s}, यानी \qty{3}{min}~\qty{56}{s} पहले होता है।

\textbf{15.} हमें चाहिए $12 + n\times24/365.2422 = 6.7525 + 24$, इसलिए
$n = 18.7525/0.065710 = 285.4$ दिन। 20 मार्च के 285 दिन बाद 30 दिसंबर है: सिरियस स्थानीय
माध्य मध्यरात्रि पर 30 से 31 दिसंबर 2026 की रात को याम्योत्तर पार करता है।

\textbf{16.} पैरानाल ग्रीनविच के \ang{70.40} पश्चिम में है: उसका स्थानीय माध्य समय
$70.40/15 = \qty{4.693}{h}$ पीछे है, यानी \qty{4}{h}~\qty{41.6}{min}।

\textbf{17.} \omterm{def:b1:spherical-astronomy:solar-time}{स्पष्ट सौर काल} = \omterm{def:b1:spherical-astronomy:solar-time}{माध्य सौर काल} + \omterm{def:b1:spherical-astronomy:solar-time}{समय का समीकरण} =
$\qty{0}{h} - \qty{3}{min}$: धूपघड़ी लगभग 23:57 दिखाती है।

\textbf{18.} $7\times\qty{3.93}{min} \approx \qty{27.5}{min}$ पहले।

\textbf{19.} जून अयनांत पर $+\ang{23.44}$, दिसंबर
अयनांत पर $-\ang{23.44}$।

\textbf{20.} जून: $\ang{90} - |\ang{-24.63} - \ang{23.44}| = \ang{41.93}$।
दिसंबर: $\ang{90} - |\ang{-24.63} + \ang{23.44}| = \ang{88.81}$; चूँकि
$\delta_\odot = \ang{-23.44} > \varphi$ है, सूर्य \omterm{def:b1:spherical-astronomy:celestial-sphere}{खमध्य} के \ang{1.19} उत्तर से गुज़रता है।

\textbf{21.} $I = 1361\times(1/0.9837)^2\times\sin\ang{88.81} = 1361\times
1.0334\times0.9998 = \qty{1406}{W.m^{-2}}$.

\textbf{22.} जून में दक्षिणी शीत ऋतु। $\cos H_0 = -\tan(\ang{-24.63})
\tan\ang{23.44} = 0.4584\times0.4335 = 0.1988$: सूर्य \omterm{def:b1:spherical-astronomy:equatorial}{घंटा कोण}
$H_0 = \ang{78.53}$ पर अस्त होता है।

\textbf{23.} $\cos\mathrm{A}_0 = \sin\ang{23.44}/\cos\ang{24.63} = 0.4376$:
$\mathrm{A}_0 = \ang{64.0}$, पूर्व के उत्तर में।

\textbf{24.} दिन की लंबाई $2\times78.53/15 = \qty{10.47}{h}$; पैरानाल पर वर्ष की सबसे लंबी रात
$24 - 10.47 = \qty{13.53}{h}$ है, यानी
\textbf{\qty[text-series-to-math=true]{13}{h}~\qty[text-series-to-math=true]{32}{min}}।
\end{solution}
